%%%%%%%%%%%%%%%%%%%%%%%%%%%%%%%%%
%
%       EXERCÍCIO
%
%%%%%%%%%%%%%%%%%%%%%%%%%%%%%%%%%

\ifdefstring{\atividade}{prova}{%
    \ifdefstring{\modo}{objetivo}{%
        \renewcommand{\valorquestao}{\ValorQObj\ ponto}
    }{%
        \renewcommand{\valorquestao}{\ValorQDisc\ pontos}
    }%
}%

\begin{exercicioBanco}[\valorquestao]
Para organizar a logística de um grande depósito, um gerente precisa estimar a quantidade de supervisores necessários. Para cada 1.000 caixas armazenadas, é necessário um supervisor.

Às 7h da manhã, o gerente observa que as caixas ocupam uma área em formato quadrado com lado de \(200\) m. Sabe-se que, em média, cada caixa ocupa \(2\ \text{m}^2\).

Nas horas seguintes, o número de caixas aumenta a uma taxa de \(30.000\) caixas por hora até as 11h, quando o depósito será fechado para novas entradas.

Determine quantos supervisores serão necessários nesse horário.

% Define as alternativas
\newcommand{\alternativas}{%
    \begin{center}
        \begin{tabularx}{\textwidth}{XXXXX}
            (a) \(130\). &
            (b) \(140\). &
            (c) \(150\). &
            (d) \(160\). &
            (e) \(170\).
        \end{tabularx}
    \end{center}
}

% Resposta correta
\newcommand{\resposta}{B}

% Lógica condicional para exibição
\ifdefstring{\atividade}{lista}{%
    \alternativas
    \vspace{0.5em}
    
    \noindent\textbf{Resposta:} letra \textbf{\resposta}.
}{%
    \ifdefstring{\modo}{objetivo}{%
        \alternativas
    }{%
        % modo = discursiva → não mostra alternativas
    }
}
\end{exercicioBanco}

